% =====================================================================
% BÁO CÁO BÀI TẬP LỚN LẬP TRÌNH NÂNG CAO -- ROOMHUB
% Biên dịch bằng pdfLaTeX; tài liệu tham khảo viết trực tiếp, không cần BibTeX.
% Thông tin bìa được đặt trong Bia.tex và Bia_lot.tex.
% Hình UML nằm trong Hinhve/*.pdf; mã nguồn TikZ của hình ở Hinhve/nguon_tikz
% (biên dịch riêng bằng XeLaTeX nếu muốn chỉnh sửa hình).
% =====================================================================
\documentclass[a4paper,12pt,oneside]{report}
\usepackage[utf8]{vietnam}
\usepackage[top=2cm, bottom=2cm, left=3.5cm, right=2.5cm]{geometry}
\usepackage{tikz}
\usetikzlibrary{arrows.meta,positioning,fit,calc}
\usepackage{graphicx} % Chèn hình ảnh
\usepackage{fancybox} % Tạo khung box
\usepackage{indentfirst} % Thụt đầu dòng ở dòng đầu tiên trong đoạn
\usepackage{amsthm} % Cho phép thêm các môi trường định nghĩa
\usepackage{latexsym} % Các kí hiệu toán học
\usepackage{amsmath} % Hỗ trợ một số biểu thức toán học
\usepackage{amssymb} % Bổ sung thêm kí hiệu về toán học
\usepackage{amsbsy} % Hỗ trợ các kí hiệu in đậm
\usepackage{times} % Chọn font Time New Romans
\usepackage{array} % Tạo bảng array
\usepackage{enumitem} % Cho phép thay đổi kí hiệu của list
\usepackage{subfiles} % Chèn các file nhỏ, giúp chia các chapter ra nhiều file hơn
\usepackage{titlesec} % Giúp chỉnh sửa các tiêu đề, đề mục như chương, phần,..
\usepackage{titletoc}
\usepackage{chngcntr} % Dùng để thiết lập lại cách đánh số caption,..
\usepackage{pdflscape} % Đưa các bảng có kích thước đặt theo chiều ngang giấy
\usepackage{afterpage}
\usepackage[ruled,vlined]{algorithm2e}  % Hỗ trợ viết các giải thuật
\usepackage{capt-of} % Cho phép sử dụng caption lớn đối với landscape page
\usepackage{multirow} % Merge cells
\usepackage{fancyhdr} % Cho phép tùy biến header và footer
% \usepackage[natbib,backend=biber,style=ieee]{biblatex} % Giúp chèn tài liệu tham khảo
\usepackage{appendix}

\usepackage[font=small,labelfont=bf,labelsep=period]{caption}

\usepackage{float}
\usepackage{subcaption}
\usepackage{xurl}

\usepackage{setspace}

% package content table
\usepackage{tocbasic}

\usepackage{tabularx} % Bảng tự co giãn cột
\usepackage{longtable} % Bảng dài nhiều trang
\usepackage{environ} % Định nghĩa môi trường bảng đặc tả use case
\usepackage{booktabs}
\usepackage{xcolor}


% ===================================================



\renewcommand\appendixname{PHỤ LỤC}
\renewcommand\appendixpagename{PHỤ LỤC}
\renewcommand\appendixtocname{PHỤ LỤC}


% Tài liệu tham khảo viết sẵn trong Tai_lieu_tham_khao.tex (không cần BibTeX)

% Môi trường bảng đặc tả use case
\NewEnviron{dacta}[2]{%
  \begin{table}[H]\centering\small\renewcommand{\arraystretch}{1.2}%
  \caption{#1}\label{#2}%
  \begin{tabularx}{\textwidth}{|>{\bfseries}p{3.0cm}|X|}\hline
  \BODY
  \hline\end{tabularx}\end{table}}
% Danh sách bước trong ô bảng: bọc minipage để chiều cao ô tính đúng, không đè dòng dưới
\newenvironment{buoc}{\begin{minipage}[t]{\linewidth}\begin{enumerate}[label=\arabic*.,leftmargin=1.4em,nosep]}{\end{enumerate}\vspace{3pt}\end{minipage}}
% Đảm bảo tên chương, hình, bảng hiển thị bằng tiếng Việt
\AtBeginDocument{%
  \renewcommand{\chaptername}{Chương}%
  \renewcommand{\figurename}{Hình}%
  \renewcommand{\tablename}{Bảng}%
  \SetAlgorithmName{Thuật toán}{thuật toán}{Danh mục thuật toán}%
}



% ===================================================


\fancypagestyle{plain}{%
\fancyhf{} % clear all header and footer fields
\fancyfoot[C]{\thepage}
\renewcommand{\headrulewidth}{0pt}
\renewcommand{\footrulewidth}{0pt}}

\setlength{\headheight}{18pt}

\def \TITLE{RoomHub -- Hệ thống đặt và quản lý phòng học}
\def \AUTHOR{Nhóm RoomHub}

% ===================================================
\definecolor{RoomHubBlue}{RGB}{20,57,108}
\titleformat{\chapter}[hang]
  {\normalfont\bfseries\color{RoomHubBlue}\fontsize{15}{19}\selectfont}
  {CHƯƠNG \thechapter.}{0.5em}{}
\titlespacing*{\chapter}{0pt}{-6pt}{14pt}
\titleformat{\section}[hang]{\normalfont\bfseries\fontsize{12}{15}\selectfont}
  {\thesection}{0.6em}{}
\titlespacing*{\section}{0pt}{12pt}{6pt}
\titleformat{\subsection}[hang]{\normalfont\bfseries}
  {\thesubsection}{0.6em}{}
\titlespacing*{\subsection}{0pt}{10pt}{4pt}
\renewcommand\thesubsubsection{\alph{subsubsection}}
\titleformat{\subsubsection}[hang]{\normalfont\bfseries}
  {\thesubsubsection)}{0.5em}{}
\titlespacing*{\subsubsection}{0pt}{8pt}{4pt}

% ===================================================
\usepackage{hyperref}
\hypersetup{pdfborder = {0 0 0}} %
\hypersetup{pdftitle={\TITLE},
	pdfauthor={\AUTHOR}}

\usepackage[all]{hypcap} % Cho phép tham chiếu chính xác đến hình ảnh và bảng biểu

\graphicspath{{Hinhve/}} % Thư mục chứa các hình ảnh

\counterwithin{figure}{chapter} % Đánh số hình ảnh kèm theo chapter. Ví dụ: Hình 1.1, 1.2,..
\counterwithin{table}{chapter}

\title{\bf \TITLE}
\author{\AUTHOR}

\setcounter{secnumdepth}{3} % Cho phép subsubsection trong report
% \setcounter{tocdepth}{3} % Chèn subsubsection vào bảng mục lục

\theoremstyle{definition}
\newtheorem{example}{Ví dụ}[chapter] % Định nghĩa môi trường ví dụ

\onehalfspacing
\raggedbottom % Tránh giãn khoảng trắng dọc giữa các đoạn
\setlength{\emergencystretch}{3em} % Giảm tràn lề với tên lớp, đường dẫn dài
%Khoảng cách xuống dòng
\setlength{\parskip}{4pt}
%Lùi đầu dòng
\setlength{\parindent}{16pt}



% =========================== BODY ===============
\begin{document}
\hypersetup{pageanchor=false}
% \newgeometry{top=2cm, bottom=2cm, left=2cm, right=2cm}
\subfile{Bia} % Phần bìa
\subfile{Bia_lot} % Phần bìa
% \restoregeometry

% ===================================================
\clearpage
\pagenumbering{arabic}
\setcounter{page}{3}
\hypersetup{pageanchor=true}
\pagestyle{plain}
\subfile{Chuong/0_2_Loi_cam_on.tex}
\clearpage
\subfile{Chuong/0_3_Tom_tat_noi_dung.tex}
\clearpage
\subfile{Chuong/0_4_Tom_tat_noi_dung_English.tex}

\clearpage
\renewcommand*\contentsname{MỤC LỤC}
\setcounter{tocdepth}{2}
\tableofcontents
\clearpage

%Tạo danh mục hình vẽ.
\renewcommand{\listfigurename}{DANH MỤC HÌNH VẼ}
{\let\oldnumberline\numberline
\renewcommand{\numberline}{Hình~\oldnumberline}
\listoffigures}
% \phantomsection\addcontentsline{toc}{section}{\numberline {} DANH MỤC HÌNH VẼ}
\newpage


 %Tạo danh mục bảng biểu.
\renewcommand{\listtablename}{DANH MỤC BẢNG BIỂU}
{\let\oldnumberline\numberline
\renewcommand{\numberline}{Bảng~\oldnumberline}
\listoftables}

\newpage
% Danh mục thuật ngữ và từ viết tắt (bảng tĩnh, không cần gói glossaries)
\chapter*{DANH MỤC THUẬT NGỮ VÀ TỪ VIẾT TẮT}
\begin{center}
\small
\renewcommand{\arraystretch}{1.15}
\begin{longtable}{|p{2.4cm}|p{11.6cm}|}
\hline
\textbf{Thuật ngữ} & \textbf{Ý nghĩa} \\ \hline
\endhead
AES & Chuẩn mã hóa nâng cao (Advanced Encryption Standard) \\ \hline
API & Giao diện lập trình ứng dụng (Application Programming Interface) \\ \hline
APK & Gói cài đặt ứng dụng Android (Android Package Kit) \\ \hline
BR & Quy tắc nghiệp vụ (Business Rule) \\ \hline
CRUD & Tạo, đọc, cập nhật, xóa (Create, Read, Update, Delete) \\ \hline
CSVC & Cơ sở vật chất \\ \hline
CTR & Chế độ mã hóa bộ đếm (Counter mode) \\ \hline
DI & Tiêm phụ thuộc (Dependency Injection) \\ \hline
DTO & Đối tượng truyền dữ liệu (Data Transfer Object) \\ \hline
HTTP & Giao thức truyền siêu văn bản (HyperText Transfer Protocol) \\ \hline
IoC & Đảo ngược điều khiển (Inversion of Control) \\ \hline
IoT & Internet vạn vật (Internet of Things) \\ \hline
JPA & Đặc tả lưu trữ Jakarta (Jakarta Persistence API) \\ \hline
JSON & Định dạng dữ liệu JavaScript Object Notation \\ \hline
JVM & Máy ảo Java (Java Virtual Machine) \\ \hline
MQTT & Giao thức nhắn tin xuất bản -- đăng ký cho thiết bị IoT (Message Queuing Telemetry Transport) \\ \hline
MVP & Sản phẩm khả dụng tối thiểu (Minimum Viable Product) \\ \hline
ORM & Ánh xạ đối tượng -- quan hệ (Object-Relational Mapping) \\ \hline
PIN & Mã số cá nhân dùng để mở khóa (Personal Identification Number) \\ \hline
QoS & Chất lượng dịch vụ (Quality of Service) \\ \hline
REST & Kiểu kiến trúc chuyển trạng thái đại diện (Representational State Transfer) \\ \hline
TLS & Giao thức bảo mật tầng giao vận (Transport Layer Security) \\ \hline
UML & Ngôn ngữ mô hình hóa thống nhất (Unified Modeling Language) \\ \hline
UUID & Định danh duy nhất toàn cục (Universally Unique Identifier) \\ \hline
\end{longtable}
\end{center}

 %\phantomsection\addcontentsline{toc}{section}{\numberline {} DANH MỤC THUẬT NGỮ VÀ TỪ VIẾT TẮT}

% \newpage
 %\subfile{chapters/0_5_Danh_muc_viet_tat.tex}

% \newpage
% \subfile{chapters/0_6_Thuat_ngu.tex}
% ===================================================


\clearpage

\pagestyle{fancy}
\fancyhf{}
\fancyhead[L]{\small\nouppercase{\leftmark}}
%\fancyhead[LE]{\rightmark}
\fancyfoot[C]{\thepage}

\chapter{GIỚI THIỆU ĐỀ TÀI}
\label{chapter:Introduction}
\subfile{Chuong/1_Gioi_thieu} % Phần mở đầu

\newpage
%\pagestyle{fancy} % Áp dụng header và footer
\chapter{KHẢO SÁT VÀ PHÂN TÍCH YÊU CẦU}
\label{chapter:Related_works}
\subfile{Chuong/2_Khao_sat}


\newpage
%\pagestyle{fancy} % Áp dụng header và footer
\chapter{NỀN TẢNG LÝ THUYẾT VÀ CÔNG NGHỆ SỬ DỤNG}
\label{chapter:Methodology}
\subfile{Chuong/3_Cong_nghe}

\newpage
%\pagestyle{fancy} % Áp dụng header và footer
\chapter{PHÂN TÍCH THIẾT KẾ, TRIỂN KHAI VÀ ĐÁNH GIÁ HỆ THỐNG}\chaptermark{THIẾT KẾ, TRIỂN KHAI VÀ ĐÁNH GIÁ}
\label{chapter:Experiment}
\subfile{Chuong/4_Ket_qua_thuc_nghiem}

\newpage
%\pagestyle{fancy} % Áp dụng header và footer
\chapter{CÁC GIẢI PHÁP VÀ ĐÓNG GÓP NỔI BẬT}
\label{chapter:SolutionAndContribution}
\subfile{Chuong/5_Giai_phap_dong_gop}
\newpage
%\pagestyle{fancy} % Áp dụng header và footer
\chapter{KẾT LUẬN VÀ HƯỚNG PHÁT TRIỂN} %Kết luận và hướng phát triển}
\label{chapter:conclusion}
\subfile{Chuong/6_Ket_luan}


% ===================================================
\newpage
\renewcommand{\bibname}{TÀI LIỆU THAM KHẢO}
\phantomsection\addcontentsline{toc}{chapter}{TÀI LIỆU THAM KHẢO}
% Danh sách tài liệu tham khảo (định dạng IEEE, viết tay để không cần chạy BibTeX)
\begin{thebibliography}{99}
\bibitem{quytrinh2026} Nhóm RoomHub, \textit{Quy trình nghiệp vụ RoomHub, phiên bản 2.1}. Tài liệu nội bộ trong thư mục \texttt{quy\_trinh\_nghiep\_vu}.
\bibitem{gamma1994design} E. Gamma, R. Helm, R. Johnson, and J. Vlissides, \textit{Design Patterns: Elements of Reusable Object-Oriented Software}. Addison-Wesley, 1994.
\bibitem{gosling2021jls} J. Gosling, B. Joy, G. Steele, G. Bracha, A. Buckley, D. Smith, and G. Bierman, \textit{The Java Language Specification, Java SE 17 Edition}. Oracle America, Inc., 2021. [Online]. Available: \url{https://docs.oracle.com/javase/specs/jls/se17/html/index.html} (visited on 09/30/2026).
\bibitem{bloch2018effective} J. Bloch, \textit{Effective Java}, 3rd ed. Addison-Wesley, 2018.
\bibitem{springboot34} Spring Team, Broadcom Inc., \textit{Spring Boot Reference Documentation, version 3.4}. [Online]. Available: \url{https://docs.spring.io/spring-boot/3.4/index.html} (visited on 09/30/2026).
\bibitem{jakartapersistence31} Eclipse Foundation, \textit{Jakarta Persistence 3.1 Specification}. [Online]. Available: \url{https://jakarta.ee/specifications/persistence/3.1/} (visited on 09/30/2026).
\bibitem{fowler2002peaa} M. Fowler, \textit{Patterns of Enterprise Application Architecture}. Addison-Wesley, 2002.
\bibitem{bernstein2009transaction} P. A. Bernstein and E. Newcomer, \textit{Principles of Transaction Processing}, 2nd ed. Morgan Kaufmann, 2009.
\bibitem{h2database} H2 Group, \textit{H2 Database Engine -- Documentation}. [Online]. Available: \url{https://www.h2database.com/html/main.html} (visited on 09/30/2026).
\bibitem{fielding2000rest} R. T. Fielding, ``Architectural styles and the design of network-based software architectures,'' Ph.D. dissertation, University of California, Irvine, USA, 2000.
\bibitem{reactnative} Meta Platforms, Inc., \textit{React Native Documentation}. [Online]. Available: \url{https://reactnative.dev/docs/getting-started} (visited on 09/30/2026).
\bibitem{typescript} Microsoft Corporation, \textit{The TypeScript Handbook}. [Online]. Available: \url{https://www.typescriptlang.org/docs/handbook/intro.html} (visited on 09/30/2026).
\bibitem{asyncstorage} React Native Community, \textit{Async Storage Documentation}. [Online]. Available: \url{https://react-native-async-storage.github.io/async-storage/} (visited on 09/30/2026).
\bibitem{mqtt5} OASIS, \textit{MQTT Version 5.0, OASIS Standard}, 2019. [Online]. Available: \url{https://docs.oasis-open.org/mqtt/mqtt/v5.0/mqtt-v5.0.html} (visited on 09/30/2026).
\bibitem{onem2mmqtt} oneM2M Partners, \textit{TS-0010 MQTT Protocol Binding}. [Online]. Available: \url{https://www.onem2m.org/technical/published-specifications} (visited on 09/30/2026).
\bibitem{dworkin2001sp80038a} M. Dworkin, ``Recommendation for block cipher modes of operation: Methods and techniques,'' National Institute of Standards and Technology, NIST Special Publication 800-38A, 2001.
\bibitem{uml251} Object Management Group, \textit{Unified Modeling Language (UML) Specification, Version 2.5.1}, 2017. [Online]. Available: \url{https://www.omg.org/spec/UML/2.5.1} (visited on 09/30/2026).
\bibitem{jest} OpenJS Foundation, \textit{Jest Documentation}. [Online]. Available: \url{https://jestjs.io/docs/getting-started} (visited on 09/30/2026).
\bibitem{martin2017clean} R. C. Martin, \textit{Clean Architecture: A Craftsman's Guide to Software Structure and Design}. Prentice Hall, 2017.
\bibitem{roomhubrepo} Nhóm RoomHub, \textit{Mã nguồn RoomHub -- classroom-booking}. [Online]. Available: \url{https://github.com/xwomen1/classroom-booking} (visited on 09/30/2026).
\bibitem{allen1983interval} J. F. Allen, ``Maintaining knowledge about temporal intervals,'' \textit{Communications of the ACM}, vol. 26, no. 11, pp. 832--843, 1983.
\end{thebibliography}


\appendices
\addappheadtotoc

% \chapter*{PHỤ LỤC} %Kết luận và hướng phát triển}

%\mainmatter
\titleformat{\chapter}[hang]{\normalfont\bfseries\color{RoomHubBlue}\fontsize{15}{19}\selectfont}{Phụ lục \thechapter.}{0.5em}{}
\titlespacing*{\chapter}{0pt}{-6pt}{14pt}

\chapter{ĐẶC TẢ USE CASE BỔ SUNG}
\label{appendix:usecase}
\subfile{Chuong/Phu_luc_A}
\newpage
\chapter{DANH SÁCH API, CÀI ĐẶT VÀ MINH CHỨNG KIỂM THỬ}\chaptermark{API, CÀI ĐẶT VÀ KIỂM THỬ}
\label{appendix:api}
\subfile{Chuong/Phu_luc_B}
\end{document}
